% !TeX root = ../papers/nc-p3.tex
\section{共通二次式の網による反対位置の対称性}
\label{sec:finite-common-net-symmetry}
\begin{lemma}\label{lem:common-net-alternating-part}
四位置を同じ四次元空間 $V$ で表し、窓を四重線型形式 $w$ とする。
四辺の関係式を対角に制限した二次式が全て同じ正則列
$Q_1,Q_2,Q_3$ の線型包に入るとする。
すると $w=G+\omega$ と一意に書ける。ここで $G$ は位置 $0,2$ の交換と
位置 $1,3$ の交換の両方で不変、$\omega\in\bigwedge^4 V^*$ である。
\end{lemma}
\begin{proof}
$V$ が生成するテンソル代数を、対称二次部分で三つの $Q_a$ の全てに
直交する元により割る。商では $x^2=\sum_a Q_a(x)z_a$ と書ける。
$[x^2,x]=0$ を $x$ の成分の三次式として比較すると各 $[z_a,x]=0$ である。
実際、正則列をなす三二次式には一次係数の関係がなく、係数空間の任意の線型汎関数を
適用できる。連続する三辺についての仮定は $w$ がこの商の次数4部分の汎関数に降りることをいう。
よって $P_x(y,z)=w(x,x,y,z)$ と置くと
\[
 w(y,x,x,z)=w(y,z,x,x)=P_x(y,z).
\]
巡回移動した窓にも同じ論法を適用して $w(x,z,y,x)=w(y,x,x,z)=P_x(y,z)$ を得る。
従って位置 $0,2$ を交換した形式と元の形式との差は、三つの連続する組のいずれに
同じベクトルを置いても零であり、完全交代形式である。位置 $1,3$ の交換についても、
四つの等式から同じことが従う。完全交代化を取ればどちらの差も
$2\operatorname{Alt}(w)$ である。$\omega=\operatorname{Alt}(w)$ が求める分解を与える。
\end{proof}

\begin{theorem}\label{thm:high-order-octad-bisymmetry}
有限な反対階数8の窓で、平面局所点の二つの定義式がともに四次以上なら、
既出の三文字または対称性判定に含まれる。
具体的に、定理\ref{thm:finite-high-order-common-octad}の共通八点切断を持つ未判定候補には
両対称性が強制され、そこで必要な非対称模型 $E_D$ に反する。
\end{theorem}
\begin{proof}
各辺の線型同定の一周の合成は八点切断を各点上スキームとして固定する。
これを $B$ と書くと $u\wedge Bu$ の成分は八点上消える二次式なので三二次式の網に入る。
$u\wedge(u\wedge Bu)=0$ と網に一次関係がないことから、$u\wedge Bu=0$ が恒等的に成り立つ。
実際 $u\wedge Bu=\sum Q_a(u)K_a$ と書けば各定数二ベクトルに $u\wedge K_a=0$ がすべての $u$ で成立する。
従って $B$ は射影的に恒等であり、四位置を整合的に同定できる。全辺の対角二次式は同じ八点を含み、
補題\ref{lem:common-net-alternating-part}を適用できる。

印を $p$ とし、位置0で $p$、位置2で $u$ を代入した位置1,3の行列を $A(u)$ とする。
この四空間は左右とも $p$ を核に含むので、$V/kp$ の三行三列の行列とみなす。
辺12の前定理の四式を対角に制限するとただ一つの二次式
$Q_{12}=sn-c r(x,x)$ を得る。$r(x,x)$ は零でない（前定理の三次恒等式と $q,r$ の独立性）。
よって $Q_{12}$ は既約で $p$ はその滑らかな点である。
四式には $A(u)$ の左核を与える三式が含まれるので、
$Q_{12}$ 上では $u$ が左核に入る。同じく辺23の二次式 $Q_{23}$ 上では $u$ が右核に入る。
いずれも $\det A(u)$ を割る。$E_D$ の式によりこの三次式には $at^2\det D$ の項があり非零である。
従って二つの既約二次因子は比例し、同じ曲面 $Q$ 上で両方の核条件が成立する。

$A(u)$ の対称部分が張る空間の次元は3以上である。実際 $A(p)$ は両隣の接平面の式の積である。
前定理の各同定が接平面を接平面に移すため、これは同じ一次式の平方の非零倍であり、
対称部分は $p$ で非零である。$Q$ は全空間を張るから、その一般点でも非零である。
$p$ からの $Q$ の射影は $\mathbf P(V/kp)$ に稠密である（一般直線は $p$ の他にもう一度交わる）。
従ってこれらの非零対称行列の射影核が平面で稠密になる。
対称行列の一本の鉛筆では不可能である。階数2の元の射影核は各一点であり、階数1の元は有限個しかない。
後者は対称階数1のVeronese曲面が直線を含まないためで、一次元以下の線型空間の場合も同じ結論となる。

反対切断を対称二乗と外積二乗に分けると、$G$ と $\omega$ は独立な対角ブロックを与える。
前段から $G$ の階数は3以上であり、非零の四形式は外積二乗に階数6の完全対合を与える。
全階数8より $\omega=0$ である。特に全ての $A(u)$ は対称になる。ところが $E_D$ の全行列を左右基底変更で
対称にできるなら、可逆な一つの行列 $J$ で $E_D J$ を対称にできる。
$E_{11}J$ の対称性より $J$ の先頭行は $(a,0,0)$ であり、$a\ne0$。
二つの $b$ 方向の対称性より $J$ の右下ブロックは $aI$、
$t$ 方向の対称性より $aD$ も対称となり、模型の $D\ne D^t$ に矛盾する。
従って未判定の場合自体が起こらない。
\end{proof}
